\begin{remark}
\label{remark-local-ring-fibre}
Let $R \to S$ be a ring map. Let $\mathfrak q$ be a prime
ideal of $S$ lying over the prime ideal $\mathfrak p$ of $R$.
According to Remark \ref{remark-fundamental-diagram}
the prime $\mathfrak q$ corresponds to a unique prime
$\overline{\mathfrak q}$ of the fibre ring
$F = S \otimes_R \kappa(\mathfrak p)$.
Then we have
$$
F_{\overline{\mathfrak q}}
\cong S_\mathfrak q \otimes_{R_\mathfrak p} \kappa(\mathfrak p)
\cong S_\mathfrak q/\mathfrak p S_\mathfrak q
$$
Namely, there is an obvious ring map
$F \to S_\mathfrak q \otimes_{R_\mathfrak p} \kappa(\mathfrak p)$
which, under the displayed isomorphism, identifies with $F \to F_{\overline{\mathfrak q}}$.
The second isomorphism follows from the fact that $\kappa(\mathfrak p)$
is the quotient of $R_\mathfrak p$ by $\mathfrak pR_\mathfrak p$.
Explicitly, write $T=R\setminus\mathfrak p$. The fibre-ring
description is $F=T^{-1}(S/\mathfrak pS)$, and the prime in it is
$\overline{\mathfrak q}=T^{-1}(\mathfrak q/\mathfrak pS)$.
For $a\in S$ and $u\in T$, the fraction
$(a+\mathfrak pS)/u$ belongs to $\overline{\mathfrak q}$
if and only if $a\in\mathfrak q$; here every $\varphi(u)$
is outside $\mathfrak q$.
Thus the mutually inverse maps between $F_{\overline{\mathfrak q}}$
and $S_{\mathfrak q}/\mathfrak pS_{\mathfrak q}$ are
$$
\frac{(a+\mathfrak pS)/u}{(b+\mathfrak pS)/v}
\longmapsto
\frac{a\varphi(v)}{\varphi(u)b}+\mathfrak pS_{\mathfrak q},
\qquad
\frac{a}{b}+\mathfrak pS_{\mathfrak q}
\longmapsto
\frac{(a+\mathfrak pS)/1}{(b+\mathfrak pS)/1},
$$
where $u,v\in T$ and $b\in S\setminus\mathfrak q$.
The first map is induced by $S\to S_{\mathfrak q}/\mathfrak pS_{\mathfrak q}$:
it kills $\mathfrak pS$, inverts $T$, and inverts exactly the
additional denominators displayed above. The second is induced by
$S\to F_{\overline{\mathfrak q}}$, which inverts
$S\setminus\mathfrak q$ and kills $\mathfrak pS$.
The formulas show that both composites fix every fraction.
The tensor comparison sends
$$
\frac{a}{b}\otimes\overline{r/t}
\longmapsto
\frac{a\varphi(r)}{b\varphi(t)}+\mathfrak pS_{\mathfrak q},
$$
for $r\in R$, $t\in T$, and has inverse
$a/b+\mathfrak pS_{\mathfrak q}\mapsto(a/b)\otimes\overline{1}$.
Balancing over $R_{\mathfrak p}$ and the vanishing of
$\mathfrak pR_{\mathfrak p}$ prove well-definedness and both
inverse identities. These maps identify the original map out of
$F$ with localization at $\overline{\mathfrak q}$.
\end{remark}